% =====================================================================
% apendice_c_dec.tex -- Apendice C: discussao sobre condicoes alem da
% base de projeto (Design Extension Conditions, DEC), adaptada da
% analise preparada para o documento de referencia deste projeto
% (claude/estrutura-nureg-1537.md), condensada e reescrita em prosa de
% apendice, e complementada com a verificacao de que o ISG-NUREG-1537 Staff
% Guidance NPR-ISG-2011-002 tambem nao aborda o tema.
% =====================================================================
\nusecbreak{C}{Ap\^endice C}
\hypertarget{apC}{}%
\section*{Ap\^endice C -- Condi\c{c}\~oes Al\'em da Base de Projeto (DEC)}
\addcontentsline{toc}{section}{Ap\^endice C -- Condi\c{c}\~oes Al\'em da Base de Projeto (DEC)}
\label{sec:apendice-c}

Este ap\^endice registra, para refer\^encia do requerente e do \'org\~ao
regulador, por que a estrutura de an\'alise de acidentes deste RPAS --
herdada da NUREG-1537 -- n\~ao inclui uma categoria formal de
\emph{Design Extension Conditions} (DEC), e discute at\'e que ponto essa
aus\^encia \'e relevante para um reator multiprop\'osito de maior porte
como o RMB.

\begin{observacao}{PTs e MTs associadas a BDB (Além da Base de Projeto)}
Neste item iremos discutir os Pareceres Técnicos da ANSN e as  Memoriais Técnicas da CNEN que trataram deste assunto em outras versões do RPAS.
\end{observacao}

\subsection*{O que a NUREG-1537 efetivamente prev\^e}

A NUREG-1537 estrutura toda a sua an\'alise de acidentes em torno do
\textbf{Acidente Hipot\'etico M\'aximo} (\emph{Maximum Hypothetical
Accident}, MHA), definido no Cap\'itulo~6 (Caracter\'isticas de Seguran\c{c}a
Projetadas) e analisado quantitativamente no Cap\'itulo~13 (An\'alise de
Acidentes) deste RPAS. Trata-se de um conceito estruturalmente distinto
do que se tornou padr\~ao para reatores de pot\^encia: um \'unico
cen\'ario, deliberadamente n\~ao-probabil\'istico, escolhido por envelope
de consequ\^encias -- a pr\'opria norma registra que, como o MHA n\~ao se
espera que ocorra, o cen\'ario n\~ao precisa ser inteiramente cre\'ivel.
N\~ao h\'a \'arvore de eventos iniciadores, an\'alise de frequ\^encia ou
espectro de cen\'arios graduados por probabilidade -- h\'a um cen\'ario
\'unico, avaliado contra os mesmos crit\'erios de dose do 10~CFR
Parte~20 e Parte~100 usados para acidentes-base de projeto comuns. Na
NUREG-1537, portanto, n\~ao existe uma categoria de eventos com crit\'erio
de aceita\c{c}\~ao mais permissivo para al\'em do MHA: ele \'e, ao mesmo
tempo, o pior cen\'ario considerado e o teto de rigor da an\'alise.

\subsection*{De onde vem o DEC, e por que n\~ao est\'a na NUREG-1537}

O conceito de \emph{Design Extension Conditions} \'e substancialmente
mais recente e nasceu em outro contexto regulat\'orio: o padr\~ao da IAEA
SSR-2/1 (\emph{Safety of Nuclear Power Plants: Design}) \citep{iaeassr21},
formalizado por volta de 2012 e refor\c{c}ado diretamente pelo acidente de Fukushima
(2011) -- quinze anos depois da publica\c{c}\~ao da NUREG-1537 (1996). O
DEC se divide em duas categorias: \textbf{DEC~A}, condi\c{c}\~oes sem
degrada\c{c}\~ao significativa do combust\'ivel (p.ex., perda de energia
externa prolongada, perda do sistema de remo\c{c}\~ao de calor residual),
e \textbf{DEC~B}, condi\c{c}\~oes com fus\~ao do n\'ucleo ou acidentes
severos (p.ex., LOCA combinado com perda total de energia, falha de
cont\^enc\~ao por \emph{bypass}). Estruturalmente, o DEC ocupa os
n\'iveis~4 e~5 da defesa em profundidade -- al\'em da base de projeto,
mas analisado com m\'etodos \emph{best-estimate} para verificar que as
libera\c{c}\~oes radiol\'ogicas permanecem dentro de limites aceit\'aveis,
tipicamente mais permissivos que os da base de projeto. A aus\^encia do
DEC na NUREG-1537 n\~ao \'e uma lacuna: \'e anacronismo. A norma \'e de
1996; o arcabou\c{c}o DEC n\~ao existia como categoria regulat\'oria
consolidada at\'e quase duas d\'ecadas depois, e nasceu especificamente
para reatores de pot\^encia.

Mesmo a atualiza\c{c}\~ao p\'os-Fukushima do pr\'oprio NRC n\~ao adotou a
terminologia DEC: a regra 10~CFR~50.155 \citep{cfr10p50155}, finalizada em 2019, trata de
mitiga\c{c}\~ao de eventos al\'em da base de projeto sob a expressão
\emph{beyond-design-basis events} (BDBE) -- n\~ao ``DEC'' --, e seu
escopo \'e explicitamente limitado a licen\c{c}as de opera\c{c}\~ao de
reatores de pot\^encia e licen\c{c}as combinadas sob a Parte~52. Reatores
n\~ao-de-pot\^encia permanecem fora do escopo mesmo na regulamenta\c{c}\~ao
americana mais recente e p\'os-Fukushima.

\subsection*{Verifica\c{c}\~ao quanto ao ISG-NUREG-1537}

O texto do NPR-ISG-2011-002 -- que complementa a NUREG-1537 para
instala\c{c}\~oes de produ\c{c}\~ao de radiois\'otopos e reatores
homog\^eneos aquosos, e que \'e a refer\^encia de ISG-NUREG-1537 adotada neste
RPAS -- foi verificado diretamente quanto a esse tema: n\~ao cont\'em
nenhuma ocorr\^encia de ``design extension'', ``DEC'' ou
``beyond-design-basis'' em sentido t\'ecnico equivalente. O ISG-NUREG-1537
amplia significativamente os Cap\'itulos~4 (Descri\c{c}\~ao do Reator),
5~(Sistemas de Refrigera\c{c}\~ao), 6~(Caracter\'isticas de Seguran\c{c}a
Projetadas), 7~(Instrumenta\c{c}\~ao), 12~(Conduta de Opera\c{c}\~oes),
13~(An\'alise de Acidentes) e 14~(Especifica\c{c}\~oes T\'ecnicas) da
NUREG-1537, mas mant\'em o MHA como o \'unico cen\'ario-teto de an\'alise
de acidentes, sem introduzir uma categoria adicional al\'em dele. A
conclus\~ao da se\c{c}\~ao anterior, portanto, permanece v\'alida tamb\'em
sob o ISG-NUREG-1537: nem a norma-base nem seu complemento vigente exigem ou
estruturam uma an\'alise de DEC para este RPAS.

\subsection*{Rela\c{c}\~ao com o quadro internacional (IAEA/WENRA)}

Fora do arcabou\c{c}o norte-americano, o quadro \'e menos uniforme. A
WENRA (\emph{Western European Nuclear Regulators Association}), no
documento \emph{``Safety Reference Levels for Existing Research
Reactors''} (novembro de 2020) \citep{wenrasrl2020}, exige explicitamente an\'alise de DEC
para reatores de pesquisa como parte da defesa em profundidade,
reproduzindo a divis\~ao DEC~A/DEC~B adaptada para instala\c{c}\~oes de
menor escala, com refer\^encia ao padr\~ao da IAEA SSR-3 (\emph{``Safety
of Research Reactors''}, 2016) \citep{iaeassr3}. Fica registrada aqui uma limita\c{c}\~ao
de verifica\c{c}\~ao: n\~ao foi poss\'ivel confirmar diretamente, a partir
do texto integral da pr\'opria SSR-3, se \'e ela que introduz a
terminologia DEC para reatores de pesquisa ou se a WENRA adapta um
conceito da SSR-3 sem reproduzir exatamente os mesmos termos -- este
ponto deve ser reverificado antes de qualquer refer\^encia normativa
direta \`a SSR-3 neste RPAS.

\subsection*{Implica\c{c}\~ao para o RMB}

A NUREG-1537, como norma-base deste Cap\'itulo~1, n\~ao pede nem
estrutura uma an\'alise de DEC -- seu arcabou\c{c}o de acidentes \'e o
MHA, sem mecanismo textual para incorporar o DEC sem reescrever partes
substanciais da l\'ogica do Cap\'itulo~13, e o ISG-NUREG-1537 n\~ao altera esse
quadro. Isso n\~ao torna o tema irrelevante para um reator multiprop\'osito
de maior porte como o RMB: h\'a evid\^encia de que reguladores europeus
(WENRA) j\'a consideram DEC aplic\'avel a reatores de pesquisa, n\~ao
apenas a reatores de pot\^encia, o que sugere que a exig\^encia \'e
tamb\'em fun\c{c}\~ao de porte e invent\'ario radiol\'ogico da
instala\c{c}\~ao, e n\~ao estritamente da categoria regulat\'oria
``pot\^encia vs. n\~ao-de-pot\^encia''. Com pot\^encia e invent\'ario de
produtos de fiss\~ao substancialmente maiores que os reatores de pesquisa
``cl\'assicos'' que a NUREG-1537 foi originalmente pensada para cobrir, o
RMB pode estar em posi\c{c}\~ao intermedi\'aria em que vale a pena avaliar,
junto ao \'org\~ao regulador nacional, se uma considera\c{c}\~ao adicional
de tipo DEC -- mesmo que n\~ao formalmente exigida pelo formato
NUREG-1537 nem pelo ISG-NUREG-1537 -- fortalece o caso de seguran\c{c}a, sem que
isso signifique abandonar o MHA como o requisito m\'inimo que a norma de
fato exige. Trata-se de uma pondera\c{c}\~ao de engenharia e de
estrat\'egia regulat\'oria a ser resolvida com o \'org\~ao regulador
nacional, n\~ao de uma conclus\~ao que os documentos-fonte permitam
extrair de forma definitiva; a discuss\~ao quantitativa do MHA em si
permanece integralmente no Cap\'itulo~13 deste RPAS.
